
\documentclass[a4paper,10pt]{article}

\usepackage[utf8]{inputenc}
\usepackage[T1]{fontenc}
\usepackage{geometry}
\usepackage{booktabs}
\usepackage{array}
\usepackage{tabularx}
\usepackage{longtable}
\usepackage{listings}
\usepackage{xcolor}
\usepackage[hidelinks]{hyperref}
\usepackage{amsmath}

% Türkçe babel paketinin "=" karakterini listings seçeneklerinde etkilemesini engeller

\geometry{
    a4paper,
    top=1.8cm,
    bottom=1.8cm,
    left=1.8cm,
    right=1.8cm
}

\setlength{\parindent}{0pt}
\setlength{\parskip}{5pt}
\renewcommand{\arraystretch}{1.18}


\lstdefinestyle{cstyle}{
    language={[ANSI]C},
    basicstyle=\ttfamily\small,
    keywordstyle=\bfseries,
    commentstyle=\itshape,
    numbers=left,
    numberstyle=\tiny,
    numbersep=7pt,
    frame=single,
    breaklines=true,
    showstringspaces=false,
    tabsize=4
}

\newcommand{\hex}[1]{\texttt{0x#1}}
\newcommand{\code}[1]{\texttt{#1}}

\title{\textbf{DEOS Network \& System Management}\\
\large Sürüm 0.3}
\author{DEOS -- Dokuz Eylül Otonom Sistemler}
\date{}

\begin{document}
\maketitle

\begin{center}
\begin{tabular}{ll}
\textbf{Doküman tipi:} & Interface Control Document (ICD) \\
\textbf{Protokol sürümü:} & v1.1 (\hex{11}) \\
\textbf{Taşıma altyapısı:} & CAN-FD ana omurga \\
\textbf{Durum:} & ICD Sürüm 0.3 -- v1.1 wire-format revizyonu \\
\end{tabular}
\end{center}


\section{SYSTEM Service Komutları}

\begin{center}
\begin{tabular}{c l}
\toprule
ID & Symbol \\
\midrule
\hex{01} & \code{DEOS\_CMD\_SYSTEM\_GET\_STATE} \\
\hex{02} & \code{DEOS\_CMD\_SYSTEM\_SET\_STATE} \\
\hex{03} & \code{DEOS\_CMD\_SYSTEM\_GET\_MODE} \\
\hex{04} & \code{DEOS\_CMD\_SYSTEM\_SET\_MODE} \\
\hex{05} & \code{DEOS\_CMD\_SYSTEM\_RESET\_NODE} \\
\hex{06} & \code{DEOS\_CMD\_SYSTEM\_GET\_NODE\_INFO} \\
\hex{07} & \code{DEOS\_CMD\_SYSTEM\_HEARTBEAT} \\
\hex{08} & \code{DEOS\_CMD\_SYSTEM\_PING} \\
\hex{09} & \code{DEOS\_CMD\_SYSTEM\_PONG} \\
\bottomrule
\end{tabular}
\end{center}

\begin{lstlisting}[style=cstyle]
typedef enum
{
    DEOS_CMD_SYSTEM_GET_STATE     = 0x01,
    DEOS_CMD_SYSTEM_SET_STATE     = 0x02,
    DEOS_CMD_SYSTEM_GET_MODE      = 0x03,
    DEOS_CMD_SYSTEM_SET_MODE      = 0x04,
    DEOS_CMD_SYSTEM_RESET_NODE    = 0x05,
    DEOS_CMD_SYSTEM_GET_NODE_INFO = 0x06,
    DEOS_CMD_SYSTEM_HEARTBEAT     = 0x07,
    DEOS_CMD_SYSTEM_PING          = 0x08,
    DEOS_CMD_SYSTEM_PONG          = 0x09
} deos_system_command_t;
\end{lstlisting}

\section{PING--PONG Ağ Erişilebilirlik Mekanizması}

DEOS protokolünde tüm adreslenebilir physical ve virtual node'lar ortak bir PING--PONG
mekanizması uygular. Bir node, erişmek istediği başka bir node'a unicast PING gönderebilir.
Geçerli PING mesajını alan hedef node, aynı Ping ID değerini içeren PONG mesajını kaynak
node'a geri göndermek zorundadır.

PING ve PONG uygulama REQUEST/RESPONSE mekanizmasından ayrı bir network-management
mekanizmasıdır. Her iki mesaj da aşağıdaki alanları kullanır:

\begin{verbatim}
Priority      = DEOS_PRIO_NETWORK
Message Class = DEOS_CLASS_NETWORK
Service       = DEOS_SERVICE_SYSTEM
\end{verbatim}

PING için \code{Command = DEOS\_CMD\_SYSTEM\_PING}; PONG için
\code{Command = DEOS\_CMD\_SYSTEM\_PONG} kullanılır.

\subsection{PING Mesaj Formatı}

PING mesajının uygulama veri alanı aşağıdaki gibi tanımlanır:

\begin{center}
\begin{tabularx}{\textwidth}{c l c X}
\toprule
Byte & Alan & Tip & Açıklama \\
\midrule
0 & Protocol Version & \code{uint8\_t} & DEOS v1.1 için \hex{11} \\
1 & Sequence & \code{uint8\_t} & Gönderen node'un TX sequence değeri \\
2 & Command & \code{uint8\_t} & \code{DEOS\_CMD\_SYSTEM\_PING} = \hex{08} \\
3 & Length & \code{uint8\_t} & \hex{04} \\
4--7 & Ping ID & \code{uint32\_t} & Gönderen tarafından üretilen correlation değeri \\
\bottomrule
\end{tabularx}
\end{center}

Toplam semantic CAN-FD data length 8 byte'tır. Ping ID, DEOS genel byte-order kuralına uygun olarak
Little-Endian taşınır. Gönderen node, aynı anda birden fazla PING işlemi yürütüyorsa aktif
işlemleri ayırt edecek Ping ID değerleri kullanmalıdır.

PING CAN Identifier alanları:

\begin{verbatim}
Priority      = DEOS_PRIO_NETWORK
Message Class = DEOS_CLASS_NETWORK
Service       = DEOS_SERVICE_SYSTEM
Destination   = ping atılan node
Source        = ping gönderen node
\end{verbatim}

\subsection{PONG Mesaj Formatı}

PONG mesajı, alınan PING mesajındaki 32-bit Ping ID değerini değiştirmeden geri taşır.

\begin{center}
\begin{tabularx}{\textwidth}{c l c X}
\toprule
Byte & Alan & Tip & Açıklama \\
\midrule
0 & Protocol Version & \code{uint8\_t} & DEOS v1.1 için \hex{11} \\
1 & Sequence & \code{uint8\_t} & PONG gönderen node'un kendi TX sequence değeri \\
2 & Command & \code{uint8\_t} & \code{DEOS\_CMD\_SYSTEM\_PONG} = \hex{09} \\
3 & Length & \code{uint8\_t} & \hex{04} \\
4--7 & Ping ID & \code{uint32\_t} & PING'den aynen echo edilen correlation değeri \\
\bottomrule
\end{tabularx}
\end{center}

PONG CAN Identifier alanları:

\begin{verbatim}
Priority      = DEOS_PRIO_NETWORK
Message Class = DEOS_CLASS_NETWORK
Service       = DEOS_SERVICE_SYSTEM
Destination   = PING mesajının Source Node değeri
Source        = PING'i alan hedef node
\end{verbatim}

PONG mesajındaki Sequence alanı PING mesajındaki Sequence değerini kopyalamaz. Her node
kendi transmit sequence sayacını kullanmaya devam eder. PING ile PONG arasındaki transaction
correlation yalnızca Ping ID üzerinden yapılır.

\subsection{Zorunlu Node Davranışı}

Geçerli ve kendisine yöneltilmiş bir PING alan her DEOS node'u aşağıdaki davranışı uygular:

\begin{enumerate}[nosep]
    \item PING mesajının Source Node ve Ping ID alanlarını okur.
    \item Yeni bir PONG mesajı oluşturur.
    \item PONG Source alanını kendi Node ID'si yapar.
    \item PONG Destination alanını PING Source değerine eşitler.
    \item PING içerisindeki Ping ID değerini değiştirmeden PONG payload'ına kopyalar.
    \item Kendi normal TX Sequence değerini kullanarak PONG'u gönderir.
\end{enumerate}

\subsection{RTT Ölçümü}

Round-trip time ölçümü gönderen node üzerinde yapılır. PING gönderildiği anda yerel zaman
kaydedilir; aynı Ping ID'yi taşıyan PONG alındığında iki yerel zaman arasındaki fark RTT
olarak hesaplanır:

\[
RTT = T_{pong} - T_{ping}
\]

Bu yöntem node'lar arasında saat senkronizasyonu gerektirmez. Timeout süresi transport veya
uygulama gereksinimine göre ayrıca tanımlanabilir; v0.3 kapsamında tek bir global timeout
süresi dondurulmamıştır.

\subsection{Broadcast PING Kuralı}

\code{DEOS\_CMD\_SYSTEM\_PING} mesajı v1.1 protokolünde yalnızca unicast Destination Node
ile kullanılmalıdır. \code{Destination = DEOS\_NODE\_BROADCAST} ile PING gönderimi geçersizdir.
Bu kural, çok sayıda node'un aynı anda PONG üretmesinden doğacak response burst trafiğini
önlemek için tanımlanmıştır.

Bir uygulama tüm node'ların erişilebilirliğini taramak istiyorsa her hedef Node ID için ayrı
unicast PING işlemi başlatmalıdır.

\subsection{HEARTBEAT Mesaj Formatı}

HEARTBEAT, bir node'un periyodik veya olay tabanlı olarak kendi yaşam bilgisini yayınlamasıdır.
PING--PONG ise belirli bir node'un başka belirli bir node'a aktif erişilebilirlik sorgusu
yapmasıdır. Bu iki mekanizma birbirinin yerine kullanılmaz.

v1.1 protokolünde HEARTBEAT mesajı boş (0 byte payload) değildir. Node'un anlık durumunu (State) taşıyan 1 byte'lık bir payload barındırır:

\begin{center}
\begin{tabularx}{\textwidth}{c l c X}
\toprule
Byte & Alan & Tip & Açıklama \\
\midrule
0 & Current State & \code{uint8\_t} & \code{deos\_state\_t} enum değeri (örn: \hex{03} ACTIVE, \hex{05} SAFE) \\
\bottomrule
\end{tabularx}
\end{center}

\end{document}
